\section{Evaluación y Selección del Stack Tecnológico}
\label{sec:stack}

\subsection{Alternativas Consideradas}
La Tabla~\ref{tab:alternativas_stack} resume las alternativas evaluadas en seis capas. Se distingue entre la tecnología adoptada para el prototipo y las opciones consideradas para una eventual ampliación. No se valoró la novedad por sí misma, sino la capacidad de aprender, integrar, probar y justificar cada herramienta.

\begin{table}[H]
    \centering
    \fontsize{8.8}{10.5}\selectfont
    \renewcommand{\arraystretch}{1.10}
    \caption{Alternativas por capa tecnológica}
    \label{tab:alternativas_stack}
    \begin{tabularx}{\textwidth}{@{}>{\bfseries\raggedright\arraybackslash}p{2.35cm} >{\raggedright\arraybackslash}p{3.65cm} >{\raggedright\arraybackslash}p{2.55cm} >{\raggedright\arraybackslash}X@{}}
        \toprule
        \textbf{Capa} & \textbf{Alternativas} & \textbf{Decisión} & \textbf{Criterio principal} \\
        \midrule
        Lenguaje principal & C++20, Rust 1.75, Python 3.12 & \textbf{Python 3.12} & Expresividad algorítmica, tipado, instrumentación y rapidez de experimentación. \\
        Red vial y grafos & OSMnx/NetworkX, Overpass API, PostGIS & \textbf{OSMnx + NetworkX} & Descarga y conversión reproducible a grafo dirigido; NetworkX se usa como oráculo. \\
        Pruebas & pytest, unittest, GoogleTest & \textbf{pytest} & Parametrización, fixtures, aserciones y compatibilidad directa con el código Python. \\
        Interfaz web & React + MapLibre, Dash, Flask + HTML & \textbf{React + FastAPI + MapLibre} & Separa la interfaz TypeScript del núcleo Python, expone contratos HTTP verificables y permite mapas interactivos sin acoplar el cálculo a la vista. \\
        Aplicación móvil & Flutter/Dart, React Native, Kivy & \textbf{Flutter 3.47 + Dart 3.13} & Comparte el mismo contrato FastAPI sin duplicar solvers; ofrece interfaz táctil adaptativa, SQLite interno del dispositivo y compilación Android verificada. \\
        Documentación & LaTeX, Typst, Word & \textbf{LaTeX + BibLaTeX/Biber} & Ecuaciones, tablas, referencias APA y control de versiones. \\
        \bottomrule
    \end{tabularx}
\end{table}

\subsection{Criterios de Comparación}
Se usaron cuatro criterios ponderados: C1 expresividad algorítmica (30\%), C2 instrumentación y aislamiento del ruido (25\%), C3 reproducibilidad y portabilidad (25\%) y C4 curva de aprendizaje (20\%). La Tabla~\ref{tab:matriz_multicriterio} presenta una escala ordinal de 1 a 5. Los puntajes son un juicio técnico del equipo, no una medición física; se incluyen para hacer explícita la decisión, no para simular precisión inexistente.

\begin{table}[H]
    \centering
    \fontsize{9.5}{11.5}\selectfont
    \renewcommand{\arraystretch}{1.15}
    \caption{Matriz multicriterio del lenguaje principal (escala 1--5)}
    \label{tab:matriz_multicriterio}
    \begin{tabularx}{\textwidth}{@{}>{\raggedright\arraybackslash}Xccccc@{}}
        \toprule
        \textbf{Alternativa} & \textbf{C1 (30\%)} & \textbf{C2 (25\%)} & \textbf{C3 (25\%)} & \textbf{C4 (20\%)} & \textbf{Ponderado} \\
        \midrule
        \textbf{Python 3.12} & 5.0 & 4.2 & 4.8 & 4.8 & \textbf{4.71} \\
        C++20 & 3.8 & 4.8 & 3.5 & 2.8 & 3.78 \\
        Rust 1.75 & 4.0 & 4.6 & 3.6 & 2.2 & 3.69 \\
        \bottomrule
    \end{tabularx}
\end{table}

\subsection{Tecnología Seleccionada y Justificación}
\begin{enumerate}
    \item \textbf{Python 3.12.8.} Permite expresar los algoritmos con tipado y pruebas sin una cadena de compilación. Las tres colas se implementaron desde cero; no se usó \texttt{heapq} para las estructuras evaluadas. La contrapartida es que los tiempos absolutos incluyen la sobrecarga del intérprete y no son extrapolables a C++ o Rust.
    \item \textbf{OSMnx 2.1.1 y NetworkX 3.7.} OSMnx descarga y convierte la red vial; NetworkX se usa como oráculo de distancias. La red se congela en JSON para que el análisis no dependa de una descarga en vivo.
    \item \textbf{pytest 9.1.1.} Pruebas parametrizadas, fixtures y ejecución selectiva. La suite completa se ejecuta con \texttt{python -m pytest python/tests -q}.
    \item \textbf{React 19.2, TypeScript 5.9, Vite 7.2, MapLibre GL 6.16 y FastAPI 0.115.} Constituyen la plataforma web del proyecto: una aplicación TypeScript publicada en GitHub Pages para consultar rutas, comparar resultados, revisar la cobertura OSM y consultar la procedencia de los datos, conectada a una API Python desplegada en Render mediante Docker que reutiliza el núcleo. Dash y Flask + HTML se consideraron alternativas, pero no se implementaron. La interfaz es un medio de interacción; el núcleo evaluado sigue siendo el algoritmo.
    \item \textbf{Flutter 3.47 y Dart 3.13 (aplicación móvil integrada).} Se implementó el cliente móvil multiplataforma en \texttt{mobile/} bajo arquitectura MVVM (\texttt{provider}, \texttt{go\_router}, \texttt{flutter\_map} y \texttt{sqflite}). Consume exactamente el mismo backend FastAPI (\texttt{v1.5.0}) para los algoritmos oficiales y la generación por semillas, evitando divergencias entre plataformas, y almacena borradores, instancias y resultados en la base SQLite interna del teléfono (sin servidor de base de datos externo). Android se compiló y verificó mediante APK de desarrollo; la compilación iOS queda estructurada para ejecutarse en un entorno macOS con Xcode.
    \item \textbf{LaTeX con BibLaTeX y Biber.} Informe reproducible, ecuaciones y referencias APA. La compilación se automatiza con \texttt{compilar.ps1}.
    \item \textbf{Servicios meteorológicos (Open-Meteo y OpenWeatherMap API).} Se implementó un diseño de doble proveedor en \texttt{src/services/weather.py}. Como opción sin credenciales se emplea Open-Meteo (datos de modelo). Como proveedor alternativo autenticado se integró \texttt{OpenWeatherMapClient} (\url{https://home.openweathermap.org/users/sign_up}). Ambas fuentes alimentan \texttt{WeatherObservation}; las consultas opcionales de la web son temporales y no reescriben el caso canónico. Una clave debe suministrarse localmente, nunca en archivos versionados.
    \item \textbf{Modelo horario de tráfico.} Reglas deterministas y coeficientes hipotéticos por franja y tipo de vía; la interfaz muestra los factores, pero aplica solo el de avenida de forma uniforme a la ruta OSM. El cliente opcional de TomTom no dispone de un experimento canónico ni de pruebas de red en este repositorio; no se declara tráfico calibrado.
\end{enumerate}

\subsection{Procedimiento de Aprendizaje e Integración}
El historial registra primero el prototipo con heap binario (02/10), luego las otras colas (04/10) y, el 28/10, los datos OSM y el banco completo. El piloto reducido se guardó el 04/11, después de esas mediciones. Estas fechas son de commits y archivos; no documentan cuándo estudió cada integrante ni acreditan sesiones de seguimiento.

\subsection{Dificultades de Adopción y Soluciones Aplicadas}
\begin{itemize}
    \item \textbf{Dependencias geoespaciales en Windows.} OSMnx depende de bibliotecas geoespaciales cuya instalación puede variar entre equipos. Las pruebas de carga usan el \textit{snapshot} JSON, sin repetir la descarga. \texttt{requirements.txt} fija versiones mínimas, no un entorno idéntico al de medición; falta comprobar una instalación desde cero.
    \item \textbf{Codificación de consola.} La consola CP1252 producía errores con caracteres Unicode. Se forzó UTF-8 en archivos y se usaron mensajes ASCII en los registros de la CLI.
    \item \textbf{Reproducibilidad de tiempos.} El banco guardado usa \texttt{perf\_counter\_ns}, calentamientos y mediana/IQR, con el recolector desactivado durante cada muestra. No se asume que una repetición produzca los mismos tiempos de pared.
    \item \textbf{Interfaz y estado.} React mantiene el estado de la sesión en el cliente; FastAPI valida cada solicitud y lee los snapshots versionados sin reescribir los JSON canónicos. Las consultas meteorológicas requieren una acción explícita y no se cachean como si fueran observaciones nuevas. La API no conserva sesiones de usuario ni resultados de ruta en el servidor.
    \item \textbf{Consulta meteorológica opcional.} Se integraron Open-Meteo y OpenWeatherMap. Cuando una clave falla o no hay red, se debe identificar como tal la respuesta guardada, sin presentarla como observación nueva. No se midieron tiempos de activación de claves.
\end{itemize}
